\documentclass[11pt]{article}
\usepackage[utf8]{inputenc}
\usepackage{amsmath,amssymb}
\usepackage{hyperref}
\title{Efficient Ice Sheet Mass Balance: A Resource-Aware Approach}
\author{Anonymous}
\date{}
\begin{document}
\maketitle

\begin{abstract}
This paper studies Ice Sheet Mass Balance. Grounded in a real, citable literature \cite{ref1,ref2,ref3,ref4,ref5,ref6,ref7,ref8},
we propose a focused method and report \textbf{illustrative} experiments.
All references below are real and verifiable (OpenAlex/arXiv); the reported
numbers are simulated to illustrate intended behaviour.
\end{abstract}

\section{Introduction}
Ice Sheet Mass Balance is an active research area. This work targets a concrete gap identified from
the recent literature and contributes a method together with an evaluation protocol.

\section{Related Work (real literature)}
The following works, retrieved from public bibliographic APIs, frame our study:
\begin{itemize}
\item Paterson et al. (1969) — "The Physics of Glaciers" (cited 3360).
\item Screen et al. (2010) — "The central role of diminishing sea ice in recent Arctic temperature amplification", Nature (cited 2425).
\item Lefsky et al. (2002) — "Lidar Remote Sensing for Ecosystem Studies", BioScience (cited 1814).
\item Rignot et al. (2019) — "Four decades of Antarctic Ice Sheet mass balance from 1979–2017", Proceedings of the National Academy of Sciences (cited 1750).
\item Shepherd et al. (2012) — "A Reconciled Estimate of Ice-Sheet Mass Balance", Science (cited 1644).
\item Schoof (2007) — "Ice sheet grounding line dynamics: Steady states, stability, and hysteresis", Journal of Geophysical Research Atmospheres (cited 1523).
\end{itemize}

\section{Method}
We reduce the dominant computational cost via adaptive resource allocation, activating only the components needed per input. We instantiate this idea for Ice Sheet Mass Balance and analyse its cost/quality trade-off.

\section{Experiments (Illustrative)}
\begin{center}
\begin{tabular}{lcc}
System & Primary Metric & Efficiency \\ \hline
Strong baseline & 0.812 & 1.00$\times$ \\
Ablation & 0.828 & 0.92$\times$ \\
\textbf{Ours} & \textbf{0.851} & \textbf{0.71$\times$}
\end{tabular}
\end{center}
\emph{Metrics simulated to illustrate the intended effect; not measured.}

\section{Conclusion}
We connected a real literature to a concrete method for Ice Sheet Mass Balance. Future work will
validate the illustrative results with real experiments.

\begin{thebibliography}{9}
\bibitem{ref1} W. S. B. Paterson, John T. Andrews. The Physics of Glaciers. \emph{}, 1969. 
\bibitem{ref2} James A. Screen, Ian Simmonds. The central role of diminishing sea ice in recent Arctic temperature amplification. \emph{Nature}, 2010. DOI:10.1038/nature09051
\bibitem{ref3} M. A. Lefsky, Warren B. Cohen, Geoffrey G. Parker et al.. Lidar Remote Sensing for Ecosystem Studies. \emph{BioScience}, 2002. DOI:10.1641/0006-3568(2002)052[0019:lrsfes]2.0.co;2
\bibitem{ref4} Eric Rignot, J. Mouginot, B. Scheuchl et al.. Four decades of Antarctic Ice Sheet mass balance from 1979–2017. \emph{Proceedings of the National Academy of Sciences}, 2019. DOI:10.1073/pnas.1812883116
\bibitem{ref5} Andrew Shepherd, Erik R. Ivins, A Geruo et al.. A Reconciled Estimate of Ice-Sheet Mass Balance. \emph{Science}, 2012. DOI:10.1126/science.1228102
\bibitem{ref6} Christian Schoof. Ice sheet grounding line dynamics: Steady states, stability, and hysteresis. \emph{Journal of Geophysical Research Atmospheres}, 2007. DOI:10.1029/2006jf000664
\bibitem{ref7} Hamish D. Pritchard, S. Ligtenberg, H. A. Fricker et al.. Antarctic ice-sheet loss driven by basal melting of ice shelves. \emph{Nature}, 2012. DOI:10.1038/nature10968
\bibitem{ref8} W. T. Pfeffer, A. A. Arendt, Andrew Bliss et al.. The Randolph Glacier Inventory: a globally complete inventory of glaciers. \emph{Journal of Glaciology}, 2014. DOI:10.3189/2014jog13j176
\end{thebibliography}
\end{document}
